\RequirePackage{fix-cm}
\documentclass[a4paper,10pt]{article}
\newcommand{\BinBookletLabel}{HEFT 2}
\newcommand{\BinPageLabel}{Arbeitsseite}
% Gemeinsames A4-Querformat mit A5-Text und karierter Notizhälfte.
\usepackage{fontspec}
\setmainfont{Latin Modern Roman}
\usepackage[ngerman]{babel}
\usepackage{amsmath,amssymb,amsthm,mathtools}
\usepackage{booktabs,tabularx,array}
\usepackage{tikz}
\usetikzlibrary{arrows.meta}
\usepackage[a4paper,landscape,top=18mm,bottom=16mm,left=12mm,right=160.5mm]{geometry}
\usepackage{enumitem}
\makeatletter
\renewcommand\normalsize{\@setfontsize\normalsize{10.5}{13}}
\makeatother
\normalsize
\setlist{nosep,leftmargin=1.5em,topsep=.35em}
\pagestyle{empty}
\usepackage{microtype}
\usepackage{needspace}
\usepackage{etoolbox}
\usepackage[hidelinks]{hyperref}
\definecolor{pascalblue}{RGB}{22,89,112}
\definecolor{pascalpale}{RGB}{231,241,244}
\newcommand{\N}{\mathbb N}
\newcommand{\GuideRef}[2]{\href{../main-results/_B10-main-results.pdf\#formula#1}{#2}}
\newcommand{\NewtonLink}{\href{../newton/_B21-newton-reading.pdf\#newton.reading}{Lesefassung zur Newtondarstellung und zu endlichen Differenzen}}
\newtheorem*{satz}{Satz}
\newtheorem*{definition}{Definition}
\theoremstyle{remark}
\newtheorem*{beispiel}{Beispiel}
\newcommand{\Beweisidee}{\par\smallskip\textbf{Beweisidee.}\ }
\newcommand{\PascalDreieck}[1]{%
\begin{tikzpicture}[x=0.62cm,y=0.52cm,
 entry/.style={minimum width=6mm,minimum height=5mm,inner sep=1pt,font=\small},
 binhighlight/.style={fill=pascalpale,text=pascalblue,rounded corners=2pt,font=\small\bfseries}]
 \foreach \n/\row in {0/{1},1/{1,1},2/{1,2,1},3/{1,3,3,1},4/{1,4,6,4,1},5/{1,5,10,10,5,1},6/{1,6,15,20,15,6,1}} {
  \node[anchor=east,font=\small,text=black!65] at (-7,-\n) {$n=\n$};
  \foreach \v [count=\k from 0] in \row {
   \ifnum#1=1
    \ifnum\n=4
     \ifnum\k=2 \node[entry,binhighlight] at ({2*\k-\n},-\n) {\v};
     \else \node[entry] at ({2*\k-\n},-\n) {\v};\fi
    \else \node[entry] at ({2*\k-\n},-\n) {\v};\fi
   \else
    \pgfmathtruncatemacro{\selected}{(\k==2 && \n>=2 && \n<=5) || (\n==6 && \k==3)}
    \ifnum\selected=1 \node[entry,binhighlight] at ({2*\k-\n},-\n) {\v};
    \else \node[entry] at ({2*\k-\n},-\n) {\v};\fi
   \fi
  }
 }
 \ifnum#1=1
  \draw[-{Stealth[length=1.5mm]},pascalblue,shorten >=3mm,shorten <=3mm] (-1,-3) -- (0,-4);
  \draw[-{Stealth[length=1.5mm]},pascalblue,shorten >=3mm,shorten <=3mm] (1,-3) -- (0,-4);
 \fi
\end{tikzpicture}}
\setlength{\parindent}{0pt}
\setlength{\parskip}{.42em plus .08em minus .08em}
\setlength{\emergencystretch}{1.5em}
\setcounter{tocdepth}{1}
\newlength{\BinWorkHeight}
\newcommand{\ExerciseSpace}[1]{%
 \par\smallskip\begingroup
 \setlength{\BinWorkHeight}{\dimexpr\pagegoal-\pagetotal-2\baselineskip\relax}%
 \ifdim\BinWorkHeight>#1mm\setlength{\BinWorkHeight}{#1mm}\fi
 \ifdim\BinWorkHeight>12mm
  \vfill{\small\sffamily Deine Rechnung oder Skizze}\par\nobreak
  \begin{tikzpicture}
   \draw[step=5mm,black!22,line width=.2pt] (0,0) grid (120mm,\BinWorkHeight);
  \end{tikzpicture}
 \fi\endgroup}
\newcommand{\Exercise}[2]{%
 \clearpage
 \ifnum#1=1
  \section{Übungsblätter}\label{bin:exercises}
  Arbeite zuerst ohne die Lösungen ab Seite~\pageref{bin:solutions}.
  Nutze das Karofeld rechts und den freien Raum unter der Aufgabe.
 \fi
 \subsection*{Aufgabe #1: #2}\phantomsection\label{bin:exercise-#1}}
\newcommand{\ExtraWork}[2]{%
 \clearpage\subsection*{Arbeitsraum zu Aufgabe #1}
 \textbf{Hinweis, falls du ihn brauchst.} #2\par
 \ExerciseSpace{130}}
\AddToHook{shipout/background}{%
 \begin{tikzpicture}[remember picture,overlay]
  \begin{scope}[shift={(current page.north west)},x=1mm,y=-1mm]
   \ifnum\value{page}>1
    \node[anchor=west,font=\sffamily\fontsize{7.5}{9}\selectfont,text=black!55]
       at (12,10) {BAND 10\quad BINOMIALKOEFFIZIENTEN\quad \BinBookletLabel};
    \draw[black!25,line width=.25pt] (12,13)--(136.5,13);
   \fi
   \node[anchor=west,font=\sffamily\fontsize{8.5}{10}\selectfont,text=black!65]
      at (161,13) {NOTIZEN\quad Fragen, Ideen und Rechnungen};
   \begin{scope}[shift={(161,22)}]
    \draw[step=5mm,black!22,line width=.2pt] (0,0) grid (125,170);
   \end{scope}
   \node[font=\sffamily\fontsize{8}{10}\selectfont,text=black!60]
      at (74.25,201) {\BinPageLabel{} \thepage};
   \node[font=\sffamily\fontsize{8}{10}\selectfont,text=black!60]
      at (223.5,201) {Notizen zu \BinPageLabel{} \thepage};
  \end{scope}
 \end{tikzpicture}}
\newcommand{\PascalEltern}{%
 \begin{tikzpicture}[x=1mm,y=1mm,
  parent/.style={draw=pascalblue!50,fill=pascalpale,rounded corners=2pt,inner sep=6pt},
  arr/.style={-{Stealth[length=2mm]},pascalblue,line width=.7pt}]
  \node[parent] (L) at (-24,0) {$\displaystyle\binom n{k-1}$};
  \node[parent] (R) at (24,0) {$\displaystyle\binom nk$};
  \node[parent] (C) at (0,-26) {$\displaystyle\binom{n+1}k$};
  \node[font=\small,anchor=east] at (-43,0) {Zeile $n$};
  \node[font=\small,anchor=east] at (-25,-26) {Zeile $n+1$};
  \draw[arr] (L.south)--(C.north west);
  \draw[arr] (R.south)--(C.north east);
  \node[font=\footnotesize,anchor=south] at (-24,8) {Index $k-1$};
  \node[font=\footnotesize,anchor=south] at (24,8) {Index $k$};
  \node[font=\footnotesize,anchor=north] at (0,-34) {Index $k$ in der neuen Zeile};
 \end{tikzpicture}}
\allowdisplaybreaks[1]
\clubpenalty=10000
\widowpenalty=10000
\BeforeBeginEnvironment{satz}{\Needspace{6\baselineskip}}
\BeforeBeginEnvironment{definition}{\Needspace{6\baselineskip}}
\BeforeBeginEnvironment{beispiel}{\Needspace{5\baselineskip}}
\pretocmd{\section}{\Needspace{9\baselineskip}}{}{}

% Leere Wegeskizzen fuer Aufgaben: keine Eintraege, Wegzahlen oder Loesungspfade.
\tikzset{
 binexpathnode/.style={circle,draw=black!45,fill=white,line width=.35pt,
   minimum size=2.2mm,inner sep=0pt},
 binexpathmark/.style={binexpathnode,draw=pascalblue,fill=pascalpale,
   line width=.7pt,minimum size=3mm},
 binexpathlabel/.style={font=\fontsize{8.5}{10}\selectfont},
 binexpathbase/.style={black!25,line width=.35pt}
}
\newcommand{\BinExercisePathEdges}[1]{%
 \pgfmathtruncatemacro{\binexlast}{#1-1}%
 \foreach \n in {0,...,\binexlast} {
  \foreach \k in {0,...,\n} {
   \draw[binexpathbase] ({2*\k-\n},-\n)--({2*\k-\n-1},-\n-1);
   \draw[binexpathbase] ({2*\k-\n},-\n)--({2*\k-\n+1},-\n-1);
  }
 }
}
\newcommand{\BinExercisePathNodes}[2]{%
 \node[binexpathlabel,anchor=east,text=black!65] at (#2,.6) {Zeile};
 \foreach \n in {0,...,#1} {
  \node[binexpathlabel,anchor=east,text=black!65] at (#2,-\n) {\n};
  \foreach \k in {0,...,\n}
   \node[binexpathnode] at ({2*\k-\n},-\n) {};
 }
 \node[binexpathmark] at (0,0) {};
 \node[binexpathlabel,anchor=south] at (0,.35) {Start};
}

\newcommand{\PascalAufgabeStart}{%
 \begin{tikzpicture}[x=7mm,y=6mm]
  \BinExercisePathEdges{4}
  \BinExercisePathNodes{4}{-5.2}
  \node[binexpathmark] at (0,-4) {};
  \node[binexpathlabel,anchor=west] at (.45,-4) {Ziel};
  \foreach \k in {0,...,4}
   \node[binexpathlabel,anchor=north] at ({2*\k-4},-4.5) {$k=\k$};
 \end{tikzpicture}%
}

\newcommand{\PascalAufgabeSchnitt}{%
 \begin{tikzpicture}[x=7mm,y=6mm]
  \fill[pascalpale,rounded corners=.8mm]
    (-5.5,-1.65) rectangle (5.5,-2.35);
  \BinExercisePathEdges{5}
  \BinExercisePathNodes{5}{-6.2}
  \foreach \x/\name in {-2/A,0/B,2/C} {
   \node[binexpathnode,draw=pascalblue!60] at (\x,-2) {};
   \node[binexpathlabel,fill=white,inner sep=.3pt] at (\x,-2.6) {\name};
  }
  \node[binexpathmark] at (1,-5) {};
  \node[binexpathlabel,anchor=west] at (1.45,-5) {Ziel};
 \end{tikzpicture}%
}

\hypersetup{hypertexnames=false,
 pdftitle={Band 10 - Binomialkoeffizienten und das Pascalsche Dreieck - Übungsheft},
 pdfauthor={Martin Kunze},
 pdfsubject={Acht Aufgaben zur Wegeidee mit kariertem Arbeitsraum und vollständigen Lösungen}}
\begin{document}
\hypertarget{binomial.exercises}{}
{\sffamily\small MATHEMATIK\quad /\quad BAND 10}\par
\vspace{3mm}\hrule\vspace{9mm}
{\fontsize{29}{32}\selectfont\bfseries Muster,\\die zählen\par}
\vspace{6mm}
{\Large Binomialkoeffizienten und\\das Pascalsche Dreieck\par}
\vspace{3mm}
{\normalsize Heft 2: Übungen und Lösungen\par}
\vspace{6mm}
\begin{center}\PascalDreieck{2}\end{center}
\vfill
{\sffamily\small Anwenden\quad\textbullet\quad Begründen\quad\textbullet\quad Nachprüfen\par}
\vspace{5mm}\hrule\vspace{3mm}
{\large Martin Kunze\par}
{\small Acht Aufgaben mit viel Raum für eigene Rechnungen\par}
\clearpage
\section*{Mit diesem Übungsheft arbeiten}
Dieses Heft ergänzt die
\href{_B10-binomial-reading.pdf\#binomial.reading}{Lesefassung
\emph{Binomialkoeffizienten und das Pascalsche Dreieck}}.
Acht Aufgaben führen vom Zeichnen einzelner Wege zum Begründen der
Formeln. Die vollständigen Lösungen beginnen auf Seite~\pageref{bin:solutions}.

\textbf{Unser Wegemodell.} Wir starten an der Spitze $(0,0)$.
Jeder Schritt führt schräg nach links unten ($L$) oder rechts unten ($R$):
\[
 L:(n,k)\longmapsto(n+1,k),\qquad
 R:(n,k)\longmapsto(n+1,k+1).
\]
Nach $n$ Schritten, davon $k$ nach rechts, stehen wir bei $(n,k)$.
Dorthin führen $\binom nk$ Wege. Der obere Index zählt alle Schritte,
der untere die Rechtsschritte. Ein Weg ohne Schritte zählt als ein Weg.

\textbf{So arbeitest du.} Zeichne oder notiere zunächst Beispiele.
Begründe dann, warum deine Einteilung jeden Weg genau einmal erfasst.
Bei einem Produkt erkläre, welche Möglichkeiten unabhängig kombiniert
werden; bei einer Summe, welche getrennten Klassen zusammenkommen.
Nutze die Karoflächen rechts und unter den Aufgaben. Für die Beweise
in Aufgabe~6 und~8 gibt es jeweils eine zusätzliche Arbeitsseite mit Hinweis.

Wie in der Lesefassung gilt $\N=\{0,1,2,\ldots\}$,
$\binom nk=0$ für $k>n$ und $a^0=1$, auch für $a=0$.
Stellen innerhalb einer Zeile zählen wir von links ab $0$.

\medskip\tableofcontents

% Jede Aufgabe erhält eine eigene Arbeitsseite; Zusatzraum bei 6 und 8.
\Exercise{1}{Wege zeichnen und die Eltern finden}
\begin{center}\PascalAufgabeStart\end{center}
\textbf{(a)} Zeichne einen Weg nach $(4,2)$ ein und notiere seine
Schrittfolge. Liste dann \emph{alle} solchen Folgen auf. Warum ist
deine Liste vollständig? Markiere jeweils die Positionen der $R$.

\textbf{(b)} Ordne die Folgen nach ihrem letzten Schritt. Von welchen
beiden Stellen kommen sie? Erkläre so $\binom42=\binom31+\binom32$.

\textbf{(c)} Übertrage die Erklärung auf $(n+1,k)$ mit $1\le k\le n+1$.
Warum heißen die Eltern $(n,k-1)$ und $(n,k)$? Was passiert bei $k=0$?

\textbf{(d)} Welche Wege zählen $\binom22,\binom32,\binom42$?
Erkläre den festen Index $2$ und die schräge Lage dieser Stellen im Dreieck.
\ExerciseSpace{40}

\Exercise{2}{Ränder, Nullen und Spiegelung}
\textbf{(a)} Bestimme die folgenden Werte und begründe jeden mit Wegen:
\[
 \binom00,\quad\binom07,\quad\binom90,\quad\binom91,
 \quad\binom98,\quad\binom99,\quad\binom9{10}.
\]

\textbf{(b)} Vertausche in jedem Weg $L$ und $R$. Wohin führt der
Spiegelweg eines Weges nach $(n,k)$? Begründe die Symmetrie durch
Hin- und Rückzuordnung, für $0\le k\le n$.

\textbf{(c)} Warum darf man $\binom35$ nicht mit
$\binom3{3-5}$ gleichsetzen? Erkläre den richtigen Wert im Wegemodell.
\ExerciseSpace{60}

\Exercise{3}{Rechtsschritte auswählen und zählen}
\textbf{(a)} Wie viele Wege nach $(12,3)$ gibt es? Erkläre zunächst,
warum ein solcher Weg durch drei ausgewählte Schrittpositionen
eindeutig bestimmt ist. Berechne die Anzahl dann mit der Fakultätenformel;
kürze vor dem Ausmultiplizieren.

\textbf{(b)} Wie viele Wege führen nach $(12,9)$? Nutze die Spiegelung.

\textbf{(c)} Beweise für jedes $n\ge2$ durch Zählen der beiden
Rechtsschrittpositionen:
\[
 2\binom n2=n(n-1).
\]
Warum zählt die rechte Seite jeden Weg genau zweimal?
Bestätige die Gleichung anschließend durch Kürzen in der Fakultätenformel.
\ExerciseSpace{55}

\Exercise{4}{Wege beim Ausmultiplizieren}
In den fünf Faktoren von $(a+b)^5$ bedeutet $R$: Wähle $a$;
$L$: Wähle $b$.

\textbf{(a)} Welchen Term liefert die Folge $RLLRL$? Wie viele
Schrittfolgen enthalten dieselbe Anzahl $R$ und $L$? Begründe damit
den Koeffizienten dieses Terms.

\textbf{(b)} Schreibe $(a+b)^5$ vollständig aus, geordnet nach
aufsteigenden Potenzen von $a$. Berechne damit $11^5$ für $a=1$, $b=10$.

\textbf{(c)} Erkläre den binomischen Lehrsatz bei $n=0$ durch den
leeren Weg und das leere Produkt. Prüfe dabei auch $a=b=0$ mit der
vereinbarten Potenzkonvention.
\ExerciseSpace{55}

\Exercise{5}{Endpunkte und gefärbte Schritte}
\textbf{(a)} Wie viele Wege mit acht Schritten gibt es insgesamt?
Ordne sie nach ihrem Endpunkt und erkläre so die Summe der Zeile $8$.

\textbf{(b)} Wie viele dieser Wege enden an einer inneren Stelle,
also weder ganz links noch ganz rechts?

\textbf{(c)} Jeder Rechtsschritt erhält nun unabhängig eine von zwei
Farben; Linksschritte bleiben ungefärbt. Zwei Wege gelten als verschieden,
wenn ihre Schrittfolge oder mindestens eine Farbe verschieden ist.
Zähle auf zwei Arten und beweise
\[
 \sum_{k=0}^n\binom nk2^k=3^n.
\]
Wie viele Färbungen besitzt ein Weg nach $(n,k)$? Wie viele Möglichkeiten
gibt es bei jedem einzelnen Schritt? Prüfe $n=0$ und bestätige das
Ergebnis mit dem binomischen Lehrsatz.
\ExerciseSpace{55}

\Exercise{6}{Nach dem letzten Rechtsschritt ordnen}
Wir betrachten alle Wege nach $(n+3,3)$, also mit drei $R$ und $n$ $L$.
Zu zeigen ist
\[
 \sum_{i=0}^n\binom{i+2}2=\binom{n+3}3.
\]
\textbf{(a)} Sei $i$ die Anzahl der $L$ \emph{vor dem letzten $R$}.
Wo endet der Wegteil vor diesem $R$? Wie viele Schritte $L$ folgen
noch? Begründe die Formel, indem du die Wege nach $i$ einteilst.
Zeige auch, wie aus jedem Anfangsweg der vollständige Weg entsteht.

\textbf{(b)} Beweise dieselbe Formel zusätzlich durch Induktion über $n$,
mit Pascal-Rekursion und Randwerten. Zitiere die Diagonalensummenformel
nicht als fertigen Satz.

\textbf{(c)} Berechne damit
\[
 1+3+6+10+15+21+28+36+45+55.
\]
\ExerciseSpace{45}
\ExtraWork{6}{Vor dem letzten $R$ liegen genau zwei $R$ und $i$ $L$;
danach sind nur noch $L$ möglich. Für die Induktion: Trenne den neu
hinzukommenden Summanden ab und suche zwei Pascal-Eltern.}

\Exercise{7}{Wege an einer Zwischenzeile trennen}
\begin{center}\PascalAufgabeSchnitt\end{center}
\textbf{(a)} Zähle für A, B und C getrennt die Hinwege von der Spitze
und die Weiterwege zum Ziel $(5,3)$. Wie viele ganze Wege gehen
jeweils über diese Station? Addiere und vergleiche mit $\binom53$.

\textbf{(b)} Erkläre, warum du innerhalb jeder Station multiplizierst
und über die Stationen addierst. Welche Rolle spielen $m=2$, $n=3$
und $r=3$ in der Vandermonde-Identität? Was zählt $\binom3{3-k}$?

\textbf{(c)} Das neue Ziel sei $(5,6)$. Erkläre mit Wegen und zusätzlich
an \emph{jedem} Summanden, warum
\[
 \sum_{k=0}^6\binom2k\binom3{6-k}=0.
\]
\ExerciseSpace{35}

\Exercise{8}{In der Mitte trennen und spiegeln}
\textbf{(a)} Trenne Wege nach $(8,4)$ nach vier Schritten.
Warum gibt es über jede Station $(4,k)$ genau $\binom4k^2$ Wege?
Berechne die Summe direkt aus Zeile $4$ und durch Zählen aller Wege zum Ziel.

\textbf{(b)} Berechne die Summe benachbarter Produkte direkt:
\[
 \sum_{k=0}^3\binom3k\binom3{k+1}.
\]
Der Faktor $\binom34$ ist $0$. Vergleiche dein Ergebnis mit $\binom64$.

\textbf{(c)} Beweise für jedes $n\in\N$:
\[
 \boxed{\displaystyle
 \sum_{k=0}^n\binom nk\binom n{k+1}=\binom{2n}{n+1}.}
\]
Trenne Wege nach $(2n,n+1)$ nach $n$ Schritten. Wähle diesmal $k$
als Anzahl der \emph{Linksschritte} im ersten Teil. Bestimme die
Rechtsschritte in beiden Teilen und nutze die Spiegelung. Erkläre
Produkt und Summe und prüfe $n=0$.
\ExerciseSpace{50}
\ExtraWork{8}{Im ersten Teil liegen $k$ Linksschritte, also $n-k$
Rechtsschritte. Wie viele der insgesamt $n+1$ Rechtsschritte fehlen
noch? Denke auch an den Fall $k=n$: Kann der zweite Teil dann gelingen?}

\clearpage
\section{Lösungen}
\label{bin:solutions}
Ein Zählbeweis erklärt, welche Wege zusammengefasst werden und warum
jeder Weg genau einmal vorkommt. Die folgenden Lösungen führen auch
die Umkehrung der Zuordnungen und die nötigen Randfälle aus.

\Needspace{7\baselineskip}
\subsection*{Zu Aufgabe 1: Wege und Eltern}
\textbf{(a)} Nach $(4,2)$ führen genau die sechs Folgen
\[
 RRLL,\quad RLRL,\quad RLLR,\quad LRRL,\quad LRLR,\quad LLRR.
\]
Ihre $R$-Positionen sind der Reihe nach
\[
 \{1,2\},\ \{1,3\},\ \{1,4\},\ \{2,3\},\ \{2,4\},\ \{3,4\}.
\]
Dies sind alle Paare aus vier Positionen. Jedes Paar legt die beiden
$R$ und damit auch die beiden $L$ eindeutig fest; kein Weg fehlt oder
kommt doppelt vor. Zum Einzeichnen gehst du zum Beispiel für $RLLR$
von $(0,0)$ über $(1,1),(2,1),(3,1)$ nach $(4,2)$.

\textbf{(b)} Auf $R$ enden $RLLR,LRLR,LLRR$: Vorher liegen drei Schritte
mit einem $R$ vor, also die Stelle $(3,1)$. Auf $L$ enden
$RRLL,RLRL,LRRL$: Sie kommen von $(3,2)$. Jede Klasse hat drei Wege;
daher $\binom42=\binom31+\binom32=3+3=6$.

\textbf{(c)} Vor dem letzten $R$ wurden erst $k-1$ der insgesamt $k$
Rechtsschritte ausgeführt; vor einem letzten $L$ schon alle $k$.
Der obere Index ist vorher jeweils $n$. Entfernen und Wiederanhängen
des letzten Schritts sind umkehrbar. Die beiden Klassen sind getrennt
und vollständig, also
\[
 \binom{n+1}k=\binom n{k-1}+\binom nk.
\]
Bei $k=n+1$ ist der zweite Beitrag $0$: In $n$ Schritten sind $n+1$
Rechtsschritte unmöglich. Bei $k=0$ gibt es nur den Weg aus lauter $L$;
seine Anzahl ist $1$, ein Index $-1$ wird nicht benötigt.

\textbf{(d)} Gezählt werden Wege mit jeweils genau zwei $R$ und
insgesamt $2$, $3$ beziehungsweise $4$ Schritten: Es sind $1$, $3$
und $6$ Wege. Bei festem $k=2$ ist jede nächste Stelle über einen
Schritt $L$ erreichbar; deshalb verläuft diese Diagonale schräg nach
links unten. Der Stellenindex gehört zur jeweiligen Zeile.

\Needspace{7\baselineskip}
\subsection*{Zu Aufgabe 2: Ränder und Spiegelung}
\textbf{(a)} Die Werte sind in der angegebenen Reihenfolge
\[
 1,\quad0,\quad1,\quad9,\quad9,\quad1,\quad0.
\]
Für $\binom00$ zählt der leere Weg. Sieben $R$ bei null Schritten
und zehn $R$ bei neun Schritten sind unmöglich. Nach $(9,0)$ führt
nur $LLLLLLLLL$, nach $(9,9)$ nur $RRRRRRRRR$. Bei $(9,1)$ kann
das einzige $R$ an jeder der neun Positionen stehen; bei $(9,8)$
gilt dasselbe für das einzige $L$.

\textbf{(b)} Aus $k$ Rechtsschritten und $n-k$ Linksschritten werden
durch Vertauschen $n-k$ Rechtsschritte und $k$ Linksschritte.
Der Spiegelweg endet bei $(n,n-k)$. Erneutes Vertauschen stellt den
ursprünglichen Weg wieder her. Somit ist die Zuordnung in beide
Richtungen eindeutig, und $\binom nk=\binom n{n-k}$.
Auch der leere Weg wird auf sich selbst abgebildet.

\textbf{(c)} Die Symmetrie setzt $0\le k\le n$ voraus. Hier ist
$5>3$, und $3-5$ ist keine zugelassene natürliche Differenz.
Fünf Rechtsschritte in nur drei Schritten sind unmöglich; also $\binom35=0$.

\Needspace{7\baselineskip}
\subsection*{Zu Aufgabe 3: Rechtsschrittpositionen}
\textbf{(a)} Wähle drei der zwölf Positionen für $R$; die übrigen sind
$L$. Umgekehrt liefert jeder Weg genau seine drei $R$-Positionen.
Die Anzahl der Wege und der Auswahlen ist deshalb dieselbe. Es gilt
\[
 3!\,9!\binom{12}3=12!=12\cdot11\cdot10\cdot9!.
\]
Weil $9!>0$ ist, dürfen wir kürzen. Mit $3!=6$ ergibt sich
$6\binom{12}3=1320$ und damit $\binom{12}3=220$.

\textbf{(b)} Die Spiegelung vertauscht drei $R$ mit neun $R$.
Daher gibt es nach $(12,9)$ ebenfalls $\binom{12}9=\binom{12}3=220$ Wege.

\textbf{(c)} Wählen wir zuerst eine Position und danach eine andere,
haben wir $n(n-1)$ Möglichkeiten. Ein Weg mit zwei $R$ entsteht
zweimal: Seine beiden $R$-Positionen können in beiden Reihenfolgen
gewählt werden. Weitere Möglichkeiten für denselben Weg gibt es nicht.
Also gilt $2\binom n2=n(n-1)$.

Die Fakultätenformel bestätigt dies:
\[
 2!\,(n-2)!\binom n2=n!=n(n-1)(n-2)!.
\]
Der gemeinsame Faktor $(n-2)!$ ist positiv, auch bei $n=2$, wo er
$0!=1$ ist. Kürzen und $2!=2$ liefern dieselbe Gleichung.

\Needspace{7\baselineskip}
\subsection*{Zu Aufgabe 4: Ausmultiplizieren als Weg}
\textbf{(a)} Die Folge $RLLRL$ liefert $abbab=a^2b^3$. Genau zwei
der fünf Faktoren müssen $a$ liefern. Das entspricht den Wegen nach
$(5,2)$, also $\binom52=10$ Möglichkeiten. Ausmultiplizieren und
Zusammenfassen ergibt deshalb den Koeffizienten $10$.

\textbf{(b)} Die Wege mit $k$ Rechtsschritten liefern jeweils
$a^kb^{5-k}$. Zeile $5$ ergibt
\[
 \begin{aligned}
 (a+b)^5={}&b^5+5ab^4+10a^2b^3\\
           &+10a^3b^2+5a^4b+a^5.
 \end{aligned}
\]
Für $a=1$, $b=10$ folgt
\[
 \begin{aligned}
 11^5&=100000+50000+10000+1000+50+1\\
     &=161051.
 \end{aligned}
\]

\textbf{(c)} Bei $n=0$ gibt es keine Faktoren zur Auswahl und genau
einen leeren Weg. Sein leeres Produkt hat den Wert $1$. Links ist
$(a+b)^0=1$, rechts steht nur $\binom00a^0b^0=1$.
Auch für $a=b=0$ gilt somit $(0+0)^0=1=\binom00\,0^0\,0^0$,
mit der vereinbarten Konvention $0^0=1$.

\Needspace{7\baselineskip}
\subsection*{Zu Aufgabe 5: Endpunkte und Farben}
\textbf{(a)} Bei jedem der acht Schritte stehen unabhängig $L$ und
$R$ zur Wahl. Daher gibt es $2^8=256$ Wege. Jeder endet an genau einer
Stelle $(8,k)$; dorthin führen $\binom8k$ Wege. Die Endpunktklassen
sind getrennt und erfassen alle Wege. Also ist
\[
 \sum_{k=0}^8\binom8k=256.
\]

\textbf{(b)} Nur die beiden Wege aus lauter $L$ oder lauter $R$ enden
am Rand. Es bleiben $256-2=254$ Wege zu inneren Stellen.

\textbf{(c)} Bei jedem Schritt wählen wir jetzt zwischen
$L$, $R_{\mathrm{rot}}$ und $R_{\mathrm{blau}}$. Das gibt insgesamt
$3^n$ gefärbte Wege. Zu einem festen ungefärbten Weg nach $(n,k)$
gibt es dagegen $2^k$ Färbungen: Jede seiner $k$ Rechtsschrittpositionen
erhält unabhängig eine der beiden Farben.

Ein gefärbter Weg bestimmt eindeutig seinen ungefärbten Weg und die
Farben an dessen Rechtsschritten. Umgekehrt ergibt jede solche Kombination
einen gefärbten Weg. Daher enthält die Klasse mit Endpunkt $(n,k)$
$\binom nk2^k$ Möglichkeiten. Addition über die getrennten, vollständigen
Endpunktklassen liefert
\[
 \sum_{k=0}^n\binom nk2^k=3^n.
\]
Bei $n=0$ besitzt der leere Weg genau eine leere Färbung:
links $\binom00\,2^0=1$, rechts $3^0=1$.
Auch der binomische Lehrsatz mit $a=2$, $b=1$ liefert unmittelbar
$(2+1)^n=\sum_{k=0}^n\binom nk2^k1^{n-k}=\sum_{k=0}^n\binom nk2^k$.

\Needspace{7\baselineskip}
\subsection*{Zu Aufgabe 6: Der letzte Rechtsschritt}
\textbf{(a)} Vor dem letzten $R$ liegen genau zwei $R$ und $i$ $L$.
Der Anfangsweg endet somit bei $(i+2,2)$; es gibt $\binom{i+2}2$
Möglichkeiten. Anschließend kommen das letzte $R$ und genau $n-i$
Schritte $L$. Dieser Rest ist eindeutig festgelegt.

Umgekehrt ergänzen wir jeden Weg nach $(i+2,2)$ durch ein $R$ und
$n-i$ Zeichen $L$. So entsteht genau ein Weg nach $(n+3,3)$, dessen
letztem $R$ genau $i$ Linksschritte vorangehen. Abschneiden und
Ergänzen sind umkehrbar. Jeder vollständige Weg gehört zu genau
einer Klasse $i\in\{0,\ldots,n\}$. Deshalb
\[
 \sum_{i=0}^n\binom{i+2}2=\binom{n+3}3.
\]

\textbf{(b) Induktionsanfang.} Bei $n=0$ steht $\binom22=\binom33$;
beide Werte sind $1$. Im Wegemodell ist dies der einzige Weg $RRR$.

\textbf{Induktionsschritt.} Die Formel gelte für ein $n\in\N$. Dann
\[
 \begin{aligned}
 \sum_{i=0}^{n+1}\binom{i+2}2
 &=\left(\sum_{i=0}^n\binom{i+2}2\right)+\binom{n+3}2\\
 &=\binom{n+3}3+\binom{n+3}2\\
 &=\binom{n+4}3.
 \end{aligned}
\]
Die zweite Gleichheit nutzt die Induktionsannahme, die dritte Pascal.
Weil $n+4=(n+1)+3$, ist dies die Behauptung für $n+1$; die Formel
gilt daher für alle $n\in\N$.

\textbf{(c)} Die zehn Zahlen sind $\binom22,\binom32,\ldots,\binom{11}2$.
Somit ist $n=9$, und die Summe beträgt $\binom{12}3=220$, wie in Aufgabe~3.

\Needspace{7\baselineskip}
\subsection*{Zu Aufgabe 7: Trennen an einer Zwischenzeile}
\textbf{(a)} Eine Station in Zeile $2$ hat die Form $(2,k)$.
Zum Ziel bleiben drei Schritte, davon $3-k$ nach rechts:
\begin{center}
\renewcommand{\arraystretch}{1.25}
\begin{tabular}{@{}ccccc@{}}
\toprule
Station & $k$ & Hinwege & Weiterwege & Ganze Wege\\
\midrule
A & $0$ & $\binom20=1$ & $\binom33=1$ & $1\cdot1=1$\\
B & $1$ & $\binom21=2$ & $\binom32=3$ & $2\cdot3=6$\\
C & $2$ & $\binom22=1$ & $\binom31=3$ & $1\cdot3=3$\\
\bottomrule
\end{tabular}
\end{center}
Insgesamt erhalten wir $1+6+3=10=\binom53$ Wege.

\textbf{(b)} An jeder Station passt jeder Hinweg zu jedem Weiterweg.
Durch Zusammenfügen entsteht ein ganzer Weg; durch Trennen nach zwei
Schritten gewinnen wir die beiden Teile eindeutig zurück. Deshalb
multiplizieren wir die Anzahlen. Jeder Weg trifft Zeile $2$ an genau
einer Station. Die drei Stationsklassen sind damit getrennt und
vollständig; deshalb addieren wir ihre Anzahlen.

Hier ist $m=2$ die Länge des ersten und $n=3$ die Länge des zweiten
Teils; insgesamt sind $r=3$ Rechtsschritte nötig. In Vandermonde lautet
die Summe $\sum_{k=0}^3\binom2k\binom3{3-k}$.
Der zusätzliche Summand $k=3$ ist $\binom23\binom30=0$.

Der Faktor $\binom3{3-k}$ zählt \emph{lokale Weiterwege} ab der
Station: drei verbleibende Schritte mit $3-k$ Rechtsschritten.
Er ist nicht die Anzahl aller Wege von der ursprünglichen Spitze
zu einem späteren Punkt.

\textbf{(c)} Sechs Rechtsschritte sind in fünf Schritten unmöglich.
Auch jeder Summand ist $0$: Für $k>2$ verschwindet $\binom2k$;
für $k\le2$ ist $6-k\ge4>3$, also $\binom3{6-k}=0$.
Damit sind sämtliche $k=0,\ldots,6$ erfasst, und die Summe ist $0$.

\Needspace{7\baselineskip}
\subsection*{Zu Aufgabe 8: Die Mitte und die Spiegelung}
\textbf{(a)} Zur Station $(4,k)$ führen $\binom4k$ Hinwege. Danach
bleiben vier Schritte mit $4-k$ Rechtsschritten, also
$\binom4{4-k}=\binom4k$ Weiterwege nach der Symmetrie.
Jeder Hinweg passt zu jedem Weiterweg; der Schnitt gewinnt beide
Teile eindeutig zurück. Der Beitrag der Station ist somit $\binom4k^2$.
Jeder Weg trifft genau eine Station, also
\[
 \sum_{k=0}^4\binom4k^2=1+16+36+16+1=70=\binom84.
\]
Die rechte Seite zählt alle Wege zum gemeinsamen Ziel $(8,4)$.

\textbf{(b)} Zeile $3$, rechts mit einer Null fortgesetzt, ist
$1,3,3,1,0$. Daher
\[
 1\cdot3+3\cdot3+3\cdot1+1\cdot0=3+9+3+0=15=\binom64.
\]

\textbf{(c)} Sei zunächst $n\ge1$. Ein Weg nach $(2n,n+1)$ hat $2n$
Schritte, darunter $n+1$ Rechtsschritte. Enthält seine erste Hälfte
$k$ Linksschritte, so enthält sie $n-k$ Rechtsschritte. Die Anzahl
der möglichen ersten Hälften ist nach der Spiegelung
\[
 \binom n{n-k}=\binom nk.
\]
In der zweiten Hälfte fehlen noch $(n+1)-(n-k)=k+1$ Rechtsschritte.
Dafür gibt es $\binom n{k+1}$ Möglichkeiten. Jeder erste Teil lässt
sich mit jedem passenden zweiten Teil verbinden; das Trennen nach
$n$ Schritten kehrt diese Verbindung um. Der Beitrag dieser Klasse
ist daher $\binom nk\binom n{k+1}$.

Jeder vollständige Weg hat eine eindeutige Zahl $k$ von Linksschritten
in seiner ersten Hälfte. Die Klassen sind getrennt und erfassen
sämtliche Wege zum Ziel. Für $k=n$ wäre die zweite Hälfte unmöglich,
weil sie $n+1$ Rechtsschritte enthalten müsste; ihr Beitrag ist $0$.
Addition über $k=0,\ldots,n$ ergibt daher
\[
 \sum_{k=0}^n\binom nk\binom n{k+1}=\binom{2n}{n+1}.
\]
Für $n=0$ ist links $\binom00\binom01=0$ und rechts $\binom01=0$.
Ein Schritt nach rechts bei insgesamt null Schritten ist unmöglich.

\end{document}
